\documentclass[conference]{IEEEtran}
\usepackage{amsmath,graphicx,booktabs,url}
\begin{document}
\title{From Flow Scores to Incident Workload: Evaluating Hybrid Intrusion Detection Across Network Captures}
\author{\IEEEauthorblockN{Author names withheld in this scaffold}}
\maketitle
\begin{abstract}
\textit{Research scaffold: complete after the reproducibility and fresh-capture experiments. State the specific finding, experimental scope, measured effect and limitations. No improvement or novelty is asserted in this draft.}
\end{abstract}
\begin{IEEEkeywords}
network intrusion detection, false alarms, evaluation, random forests
\end{IEEEkeywords}
\section{Introduction}
State the operational problem and research question. List contributions only after they have been established by the experiments.
\section{Related Work}
Random Forest and hybrid detection have established precedents~\cite{zhang2008}. Cross-evaluation and generalization are also established research areas~\cite{apruzzese2022,cantone2024}. Discuss evaluation pitfalls~\cite{arp2022} and the overlap with direct-PCAP pipelines~\cite{cho2026} and calibrated alerting under benchmark shift~\cite{jin2026}.
\section{System and Experimental Protocol}
Describe the shared 22-feature extractor, five-second windows, detector semantics, frozen capture groups and validation-only operating-point selection. Distinguish intrusion candidates from informational or policy-related events. Define all denominators and coverage exclusions.
\section{Results}
\textit{Pending measured evidence. Insert reviewed tables generated by the research runner; identify published-test diagnosis and prospective evaluation separately. Do not interpret a smoke-run score as benchmark evidence.}
\section{Discussion and Limitations}
Address capture dependence, label coverage, class prevalence, detector coverage, fixed-window grouping and the distinction between connection recall and attack-episode recall. State that offline inference measurements do not establish end-to-end latency. Explain which conclusions require more independent environments or a second dataset.
\section{Conclusion}
\textit{Complete after assessing the scientific contribution.}
\bibliographystyle{IEEEtran}
\bibliography{references}
\end{document}

